\lecture{40}{Децентрализованный SGD}

\sourcevideo
    {https://www.youtube.com/playlist?list=PLOzRYVm0a65fhKDS8cYIC7YCuVGJ4FOJx}
    {Week 11: Lecture 40: Decentralized Stochastic Gradient Descent}

В параллельном обучении все работники используют одну модель и после
каждого шага глобально усредняют градиенты. Децентрализованный SGD
отказывается от центрального сервера: агент хранит собственную модель,
обрабатывает локальные данные и обменивается параметрами только с
соседями. Поэтому алгоритм должен одновременно решать задачу
оптимизации и подавлять рассогласование локальных моделей.

\section{Задача и параллельный SGD}

Пусть агент \(i\) хранит распределение или набор данных
\(\mathcal D_i\) и минимизирует локальную функцию
\[
    f_i(x)
    =
    \mathbb E_{\xi\sim\mathcal D_i}
    \bigl[\ell(x;\xi)\bigr].
\]
Глобальная задача имеет вид
\[
    \min_{x\in\mathbb R^d}F(x),
    \qquad
    F(x)=\frac1N\sum_{i=1}^{N}f_i(x).
\]
Её децентрализованная форма записывается как
\[
\begin{aligned}
    \min_{x_1,\ldots,x_N}\quad
    &\frac1N\sum_{i=1}^{N}f_i(x_i),
    \\
    \text{при условии}\quad
    &x_1=\cdots=x_N.
\end{aligned}
\]
Каждый агент знает только свою функцию и не обязан передавать
исходные данные другим участникам.

На одном узле SGD записывается как
\[
    x^{k+1}=x^k-\eta_k g(x^k;\xi_k),
    \qquad
    \mathbb E[g(x;\xi)]=\nabla F(x).
\]
При несмещённости и ограниченной дисперсии градиента типичная оценка
стационарности для гладкой функции содержит стохастический член
порядка
\[
    \frac1T\sum_{k=0}^{T-1}
    \mathbb E\lVert\nabla F(x^k)\rVert^2
    =
    O\!\left(
        \frac{\sigma}{\sqrt T}
    \right),
\]
наряду с начальным функциональным разрывом и константой гладкости.

В параллельном SGD все \(N\) работников используют одну модель
\(x^k\), независимо вычисляют
\[
    g_i^k=g_i(x^k;\xi_i^k)
\]
и выполняют глобальное обновление
\[
    x^{k+1}
    =
    x^k-
    \eta_k\frac1N\sum_{i=1}^{N}g_i^k.
\]
Если шумы независимы и имеют дисперсию не более \(\sigma^2\), то
дисперсия среднего градиента уменьшается примерно до
\(\sigma^2/N\), а стохастический член принимает порядок
\[
    O\!\left(
        \frac{\sigma}{\sqrt{NT}}
    \right).
\]
Это линейное ускорение по числу работников требует достаточной
независимости шумов и пренебрежимо малой стоимости синхронизации.

\section{Коммуникационный граф и алгоритм}

Глобальное усреднение поддерживает равенство всех моделей после
каждой итерации, но создаёт узкое место. У сервера параметров нагрузка
центрального узла имеет порядок \(O(Nd)\), а ring all-reduce требует
\(2(N-1)\) последовательных коммуникационных стадий. В
децентрализованной схеме обмен задаётся связным неориентированным
графом
\[
    \mathcal G=(\mathcal V,\mathcal E).
\]
Агент \(i\) взаимодействует только с вершинами
\(j\in\mathcal N_i\). Матрица смешивания
\(W=(w_{ij})\in\mathbb R^{N\times N}\) удовлетворяет
\[
    w_{ij}\ge0,
    \qquad
    w_{ij}=0
    \quad
    \text{при }
    j\notin\mathcal N_i\cup\{i\},
\]
\[
    W\mathbf1=\mathbf1,
    \qquad
    \mathbf1^{\mathsf T}W=\mathbf1^{\mathsf T}.
\]
Таким образом, \(W\) является дважды стохастической.

На итерации \(k\) агент вычисляет локальный стохастический градиент
\[
    g_i^k=g_i(x_i^k;\xi_i^k),
\]
делает локальный шаг
\[
    x_i^{k+\frac12}
    =
    x_i^k-\eta_k g_i^k
\]
и смешивает промежуточные модели соседей:
\[
    x_i^{k+1}
    =
    \sum_{j=1}^{N}
    w_{ij}x_j^{k+\frac12}.
\]
Итоговая рекурсия равна
\[
    x_i^{k+1}
    =
    \sum_{j=1}^{N}
    w_{ij}
    \bigl(
        x_j^k-\eta_k g_j^k
    \bigr).
\]
После объединения локальных векторов
\[
    \mathbf x^k
    =
    \begin{pmatrix}
        x_1^k\\
        \vdots\\
        x_N^k
    \end{pmatrix},
    \qquad
    \mathbf g^k
    =
    \begin{pmatrix}
        g_1^k\\
        \vdots\\
        g_N^k
    \end{pmatrix},
\]
получаем матричную форму
\[
    \mathbf x^{k+1}
    =
    (W\otimes I_d)
    \bigl(
        \mathbf x^k-\eta_k\mathbf g^k
    \bigr).
\]
Можно сначала смешать модели, а затем вычислить локальный градиентный
шаг. Обе схемы используют те же два механизма, но порождают разные
траектории; далее рассматривается порядок «локальный шаг, затем
смешивание».

\section{Средняя модель и рассогласование}

Определим среднюю модель
\[
    \overline x^k
    =
    \frac1N
    \sum_{i=1}^{N}x_i^k.
\]
Благодаря столбцовой стохастичности
\[
\begin{aligned}
    \overline x^{k+1}
    &=
    \frac1N
    \sum_{i,j}
    w_{ij}
    \bigl(
        x_j^k-\eta_k g_j^k
    \bigr)
    \\
    &=
    \overline x^k
    -
    \frac{\eta_k}{N}
    \sum_{j=1}^{N}g_j^k.
\end{aligned}
\]
Следовательно,
\[
    \overline x^{k+1}
    =
    \overline x^k-\eta_k\overline g^k,
    \qquad
    \overline g^k
    =
    \frac1N
    \sum_{i=1}^{N}g_i^k.
\]
Это похоже на параллельный SGD, но градиенты вычисляются в различных
точках \(x_i^k\). Поэтому \(\overline g^k\) не является точным
стохастическим градиентом в точке \(\overline x^k\).

Пусть
\[
    J
    =
    \frac1N
    \mathbf1\mathbf1^{\mathsf T}
\]
и
\[
    \mathbf e^k
    =
    \mathbf x^k
    -
    \mathbf1\otimes\overline x^k.
\]
Чистый смешивающий шаг действует на неконсенсусную компоненту
матрицей \(W-J\). Введём коэффициент смешивания
\[
    \rho
    =
    \left\lVert
        W-
        \frac1N
        \mathbf1\mathbf1^{\mathsf T}
    \right\rVert_2.
\]
Для связной апериодической сети
\[
    0\le\rho<1,
\]
и без градиентного возмущения
\[
    \lVert\mathbf e^{k+1}\rVert
    \le
    \rho\lVert\mathbf e^k\rVert.
\]
Обычно спектральным разрывом называют \(1-\rho\), а не саму
величину \(\rho\). Для симметричной дважды стохастической матрицы
\[
    \rho
    =
    \max_{r\ge2}
    |\lambda_r(W)|.
\]

\section{Топология и скорость смешивания}

Для полного графа можно выбрать
\[
    W
    =
    \frac1N
    \mathbf1\mathbf1^{\mathsf T},
\]
поэтому \(\rho=0\) и точное среднее достигается за один раунд. Цена
состоит в обмене каждого агента с \(N-1\) соседями.

В пути и кольце степень вершины ограничена константой, поэтому один
раунд дешёв, но информация распространяется медленно. Для стандартных
локальных весов на кольце
\[
    \rho
    =
    1-
    \Theta\!\left(
        \frac1{N^2}
    \right),
\]
следовательно,
\[
    1-\rho
    =
    \Theta\!\left(
        \frac1{N^2}
    \right).
\]
Число раундов, необходимое для существенного уменьшения
рассогласования, поэтому растёт квадратично по порядку \(N\).

Диаметр графа равен \(N-1\) для пути,
\(\lfloor N/2\rfloor\) для кольца и единице для полного графа.
Малый диаметр обычно помогает распространению информации, но не
определяет скорость смешивания полностью: точная характеристика
задаётся спектром \(W\).

При выборе сети приходится одновременно контролировать максимальную
степень, диаметр и коэффициент \(\rho\). Полный граф минимизирует
число раундов, а путь и кольцо минимизируют локальную степень.
Экспоненциальные графы дают промежуточный вариант: небольшую степень и
существенно более быстрое распространение информации, чем в кольце.

\section{Ошибка и коммуникационная стоимость}

При липшицевости локальных градиентов
\[
    \left\lVert
        \nabla f_i(x_i^k)
        -
        \nabla f_i(\overline x^k)
    \right\rVert
    \le
    L_i
    \lVert
        x_i^k-\overline x^k
    \rVert.
\]
Поэтому рассогласование моделей непосредственно искажает направление
среднего градиентного шага. Анализ децентрализованного SGD должен
одновременно контролировать оптимизационную ошибку средней модели,
стохастический шум и ошибку консенсуса. Первая уменьшается за счёт
градиентных шагов, вторая зависит от дисперсии и размеров
мини-пакетов, а третья определяется графом и матрицей смешивания.

На одной итерации агент \(i\) передаёт вектор размерности \(d\)
каждому соседу. Его локальная коммуникационная нагрузка имеет порядок
\[
    O(d_i d),
\]
а максимальная нагрузка по сети --- порядок
\[
    O(d_{\max}d),
    \qquad
    d_{\max}
    =
    \max_i d_i.
\]
Разреженный граф уменьшает стоимость одного раунда, но обычно
увеличивает число итераций из-за медленного смешивания. Поэтому
сравнивать архитектуры следует по полной стоимости достижения заданной
точности, а не только по числу передаваемых компонент за один шаг.

\section{Предпосылки и границы модели}

Базовый анализ предполагает синхронные итерации, фиксированный
связный граф, дважды стохастическую матрицу смешивания, точную передачу
параметров и несмещённые локальные градиенты. Линейное ускорение
параллельного SGD дополнительно требует независимости шумов. При
неодинаковых распределениях данных, задержках, потерях сообщений,
отказах и асинхронности возникают дополнительные ошибки и нужны
модифицированные алгоритмы.

Децентрализация устраняет центральный узел и глобальную
синхронизацию, но не делает коммуникацию бесплатной. Итоговый темп
определяется совместным действием стохастического шума, локальной
геометрии задачи и спектральных свойств сети.
